%!TEX root = ../thesis.tex

\chapter{EfficientZero}\label{chap:efficientzero}

EfficientZero~\cite{ye_mastering_2021} is a model-based deep reinforcement
learning algorithm and a member of the AlphaGo family of algorithms.
It combines a learned latent-space model of the environment with Monte-Carlo
tree search to simulate the outcome of a sequence of actions.
EfficientZero was the first algorithm to exceed average human performance on
Atari games with only two hours of real-time game experience, reaching
\SI{194.3}{\percent} mean human performance on the Atari 100k
benchmark~\cite{ye_mastering_2021}.

\section{From AlphaGo to EfficientZero}\label{sec:ez-lineage}

AlphaGo~\cite{silver_mastering_2016} was released in 2016 and was the first program to defeat a
professional human player at the game of Go, a challenging task due to Go's extremely large state space.
It combines Monte-Carlo tree search with two deep neural networks: a policy
network that proposes promising moves and narrows the search, and a value
network that evaluates board positions.
The networks were first trained by supervised learning on a database of human
expert games and then refined through self-play.
AlphaGo is therefore tied to Go since it requires human game records for
pretraining and a perfect simulator of the game rules.

AlphaGo Zero~\cite{silver_mastering_2017-1} removes the dependency on human
data.
It learns exclusively from self-play, using a single network with a policy and a value head.
The training loop is an instance of GPI (\cref{sec:rl-dp}): MCTS improves
the policy and the improved policy is used to generate new games to train the
networks to predict the output of MCTS.
AlphaGo Zero defeated the published AlphaGo by 100 games to 0.

AlphaZero~\cite{silver_general_2018} rids itself entirely of the Go-specific dependencies.
The resulting algorithm is general enough to be trained not only on Go, but
also on other board games like chess and shogi and defeated the strongest
existing programs in each of the three games.
AlphaZero still requires a perfect simulator for MCTS.

MuZero~\cite{schrittwieser_mastering_2020} removes the dependency on the simulator.
Instead of planning with the true environment dynamics, MuZero learns a model
of the environment and plans inside it.
Unlike previous approaches, the model is not trained to predict actual observations,
but encodes observations through a representation network and plans in a latent space.
This makes the approach applicable to domains whose rules are unknown, and
MuZero matched AlphaZero's performance in board games while simultaneously
setting a new state of the art on the Atari benchmark.

EfficientZero~\cite{ye_mastering_2021} improves the sample efficiency of MuZero.
It targets the Atari 100k benchmark~\cite{kaiser_model-based_2024}, a benchmark
set that wraps many classic Atari games and allows an agent to interact with each
game for only \num{100000} steps which is equal to roughly two hours of real-time play.
It achieves this through three additions to MuZero, described in
\cref{sec:ez-improvements}.
A successor, EfficientZero V2~\cite{wang_efficientzero_2024}, extends the
approach to continuous action spaces and further benchmarks.

\section{General Architecture}\label{sec:ez-architecture}

The core architecture of EfficientZero is mostly similar to MuZero.
Its world model consists of a representation network $h$, a dynamics network $g$,
a value network $v$, a value-prefix network $\vpnet$, and a policy network $p$,
wired together as shown in \cref{fig:ez-architecture}.

%!TEX root = ../thesis.tex
% Architecture overview for \cref{sec:ez-architecture}: the five networks of
% EfficientZero and the flow between them, unrolled two imagined steps.
% Superscripts count imagined steps, as in the surrounding prose; subscripts are
% reserved for real time steps ($o_t$), following MuZero's convention.
\begin{figure}[htbp]
  \centering
  \begin{tikzpicture}[
      >={Stealth[round]}, font=\small,
      net/.style    = {draw, rounded corners=2pt, fill=blue!8,
                       minimum height=7.5mm, minimum width=9mm},
      latent/.style = {draw, circle, fill=black!8, minimum size=9.5mm,
                       inner sep=0pt},
      sig/.style    = {inner sep=1.5pt},
      flow/.style   = {->, thick, black!70},
      carry/.style  = {->, thick, black!70, densely dashed},
    ]
    % --- representation network and the latent unroll ------------------------
    \node[sig]    (o)  at (-0.3,0) {$o_1,\dots,o_t$};
    \node[net]    (R)  at ( 1.6,0) {$h$};
    \node[latent] (s0) at ( 3.4,0) {$s^0$};
    \node[net]    (G1) at ( 5.2,0) {$g$};
    \node[latent] (s1) at ( 7.0,0) {$s^1$};
    \node[net]    (G2) at ( 8.8,0) {$g$};
    \node[latent] (s2) at (10.6,0) {$s^2$};
    \node[sig]    (sK) at (12.1,0) {$\cdots$};

    \draw[flow] (o)  -- (R);   \draw[flow] (R)  -- (s0);
    \draw[flow] (s0) -- (G1);  \draw[flow] (G1) -- (s1);
    \draw[flow] (s1) -- (G2);  \draw[flow] (G2) -- (s2);
    \draw[flow] (s2) -- (sK);

    % --- actions entering the dynamics network -------------------------------
    \node[sig] (a1) at (5.2,-1.3) {$a^1$};
    \node[sig] (a2) at (8.8,-1.3) {$a^2$};
    \draw[flow] (a1) -- (G1);
    \draw[flow] (a2) -- (G2);

    % --- policy and value heads, evaluated at every latent state -------------
    \node[net] (P0) at ( 2.6,1.5) {$p$};  \node[net] (V0) at ( 4.2,1.5) {$v$};
    \node[net] (P1) at ( 6.2,1.5) {$p$};  \node[net] (V1) at ( 7.8,1.5) {$v$};
    \node[net] (P2) at ( 9.8,1.5) {$p$};  \node[net] (V2) at (11.4,1.5) {$v$};

    \draw[flow] (s0) -- (P0);  \draw[flow] (s0) -- (V0);
    \draw[flow] (s1) -- (P1);  \draw[flow] (s1) -- (V1);
    \draw[flow] (s2) -- (P2);  \draw[flow] (s2) -- (V2);

    \draw[flow] (P0) -- ++(0,0.7) node[sig,above] {$p^0$};
    \draw[flow] (V0) -- ++(0,0.7) node[sig,above] {$v^0$};
    \draw[flow] (P1) -- ++(0,0.7) node[sig,above] {$p^1$};
    \draw[flow] (V1) -- ++(0,0.7) node[sig,above] {$v^1$};
    \draw[flow] (P2) -- ++(0,0.7) node[sig,above] {$p^2$};
    \draw[flow] (V2) -- ++(0,0.7) node[sig,above] {$v^2$};

    % --- recurrent value-prefix network over the imagined latents -----------
    \node[net] (VP1) at ( 7.0,-2.5) {$\vpnet$};
    \node[net] (VP2) at (10.6,-2.5) {$\vpnet$};

    \draw[flow] (s1) -- (VP1);
    \draw[flow] (s2) -- (VP2);
    \draw[carry] (VP1) -- node[above,inner sep=3pt,font=\scriptsize] {LSTM state} (VP2);
    \draw[carry] (VP2) -- ++(1.1,0) node[sig,right] {$\cdots$};

    \draw[flow] (VP1) -- ++(0,-0.7) node[sig,below] {$\hat{x}^1$};
    \draw[flow] (VP2) -- ++(0,-0.7) node[sig,below] {$\hat{x}^2$};
  \end{tikzpicture}
  \caption[The EfficientZero architecture]{%
    The five networks of EfficientZero, unrolled for two imagined steps.
    The representation network $h$ encodes the observation history into the root
    latent state $s^0$.
    The dynamics network $g$ advances a latent state by one action,
    $s^k = g(s^{k-1},a^k;\theta)$.
    Every latent state is scored by the policy head $p$ and the value head $v$.
    The value-prefix network $\vpnet$ is recurrent: given a sequence of
    latent states, it predicts the accumulated reward (\cref{sec:ez-value-prefix}).
  }
  \label{fig:ez-architecture}
\end{figure}


\paragraph{Representation network $h$}
The representation network $h:\mathcal{O}\to\mathcal{S}$ maps observations $o_t$
to a latent representation $s_t$.
Other networks operate only on latent representations, so it can be interpreted
as extracting information from the observations that is relevant for achieving
a high reward, while discarding non-relevant information~\cite{grimm_value_2020}.
This sets it apart from world models that are trained to predict or reconstruct
raw observations~\cite{ha_world_2018,kaiser_model-based_2024,hafner_learning_2019,hafner_mastering_2024},
where the model must also capture detail that is irrelevant to the task.

\paragraph{Dynamics network $g$}
The dynamics network $g:\mathcal{S}\times\mathcal{A}\to\mathcal{S}$ maps a
current latent $s^i$ and an action $a^i$ to the next latent
$s^{i+1}=g(s^i,a^i;\theta)$.
A single forwarding step is known as a \emph{rollout} step.
$g$ can be interpreted as a learned model of the environment behavior.
Instead of actually taking action $a^i$ and then re-encoding the resulting
observation as $h(o^{i+1};\theta)$, we can \emph{imagine} forwarding the latent.
This is a key part of the high sample efficiency of EfficientZero.
Throughout this chapter, $s^0$ denotes the encoded root latent and $s^k$ for
$k>0$ the latent reached after $k$ imagined steps; where the encoded latent at
a later real time step is meant, it is written out as $h(o_{t+k};\theta)$.

\paragraph{Value network $v$}
The value network $v:\mathcal{S}\to\mathbb{R}$ estimates the expected discounted
return (\cref{eq:value-function}) of a state $s_t$.

\paragraph{Value-prefix network $\vpnet$}
The value prefix $\vpnet:\mathcal{S}\times\mathcal{A}^k\to\mathbb{R}$ estimates
the expected discounted sum of all rewards when starting in a state $s_t$ and
applying actions $(a^1,\dots,a^k)$.
Combining the value network $v$ with the value-prefix network $\vpnet$ allows
MCTS to guide its search towards high-valued actions and states.

\paragraph{Policy network $p$}
The policy network $p:\mathcal{S}\to\Delta(\mathcal{A})$ maps a state $s_t$ to a
probability distribution over actions, whose component for action $a$ is written
$p_a(s)$.
It models a learned policy $\pi(a|s)$ and acts as a search prior for MCTS.

\section{Monte-Carlo Tree Search}\label{sec:ez-mcts}

Monte-Carlo tree search
(MCTS)~\cite{hutchison_bandit_2006,hutchison_efficient_2007} is a tree search
algorithm. It makes use of the UCT rule that derives from the UCB rule (\cref{eq:ucb}) introduced in
\cref{sec:rl-exploration}.
It searches through the action space $\mathcal{A}$ to find an action that maximizes the
expected future return.
On a high level, the search is guided by the policy network $p$, which acts as
a prior distribution over the possible actions, as well as the value network
$v$ and value-prefix network $\vpnet$ to prefer higher-valued states.
The use of a prior entails a slight modification of the UCT rule, called pUCT (prior UCT).
The dynamics network $g$ is used to \enquote{imagine} the successor states when
taking certain actions without actually needing to interact with the environment.
MCTS is applied in each state and produces a root visit count distribution as well
as a root value estimate.
The root visit count distribution can then be sampled to choose an action with
which to step the environment.

MCTS itself is a general framework that can be adapted to many situations.
It can be interpreted as a policy improvement operator: given a policy prior
and search heuristics, it produces a new, better policy.
Previous algorithms like AlphaGo did not use imagined rollouts, but relied on a perfect simulator
of the Go rules.
The use of imagined rollouts is a key factor in the sample efficiency of EfficientZero.

The MCTS variant used in EfficientZero consists of three phases that are
repeated for a fixed number of times: \emph{selection}, \emph{expansion}, and
\emph{backup}.
MCTS stores the current state of its search tree in nodes, where
already visited nodes are marked as \emph{expanded}.
Each node stores a latent state $s^k$, its visit count $N$, and the sum of the
returns $Q_\Sigma$ backed up through it.
\Cref{fig:ez-mcts-simulation} traces a single simulation through the three
phases on a worked example.

%!TEX root = ../thesis.tex
% One MCTS simulation for \cref{sec:ez-mcts}, traced through the three phases.
% Worked example: |A| = 2 and gamma = 1 for legibility. The tree's Q values span
% exactly [0,1], so the tree-wide normalization of \cref{eq:ez-qnorm} is the
% identity here and the displayed values are both Q and \bar{Q}.
%
% Layout: the three phases are stacked as full-width rows rather than set side by
% side, and every tree grows LEFT TO RIGHT. Three 0.32\textwidth columns left only
% ~4.6cm per panel, which forced \tiny labels; one full-width row per phase gives
% ~8.7cm of drawing width and lets the statistics be set in \scriptsize. The
% subcaption sits beside its tree, not underneath it, so the rows stay short.
% Layout rule, to keep the panels collision-free: the simulation path runs
% straight right along y=0 with siblings on the row below; statistics of path
% nodes are therefore always set ABOVE, statistics of siblings always BELOW, and
% edge labels sit on the edge itself with a white background.
% Every panel opens with the same invisible \path at y=0, which pads all three
% tikzpictures to the same horizontal extent so s^0, s^1, s^2 line up in a column
% across the rows. It deliberately does not fix the vertical extent: (b) and (c)
% carry no statistics under their sibling nodes and would otherwise reserve the
% empty space that only (a) needs.
\begin{figure}[htbp]
  \centering
  % The subcaptions sit in a 0.36\textwidth column, too narrow to justify without
  % gaping word spaces -- set them ragged right for this figure only.
  \captionsetup[sub]{justification=raggedright, singlelinecheck=false}
  \tikzset{
    >={Stealth[round]},
    nd/.style      = {draw, circle, fill=black!8, minimum size=8mm,
                      inner sep=0pt, font=\small},
    ndpath/.style  = {nd, fill=blue!15, very thick},
    ndnew/.style   = {nd, fill=blue!15, very thick, draw=blue!60!black},
    ndopen/.style  = {draw, circle, densely dashed, black!50, minimum size=8mm,
                      inner sep=0pt, font=\small},
    net/.style     = {draw, rounded corners=2pt, fill=white, font=\scriptsize,
                      inner sep=2pt},
    ed/.style      = {->, black!55},
    edsel/.style   = {->, very thick, blue!60!black},
    edup/.style    = {->, very thick, densely dashed, black!75},
    pr/.style      = {font=\scriptsize, inner sep=1pt, black!55, fill=white},
    prsel/.style   = {font=\scriptsize, inner sep=1pt, fill=white},
    st/.style      = {font=\scriptsize, inner sep=1pt, align=center},
  }

  % ===== (a) Selection =====================================================
  \begin{subfigure}[b]{\textwidth}
    \begin{minipage}[c]{0.60\textwidth}
      \centering
      \begin{tikzpicture}
        \path (-1.1,0) (7.5,0);
        \node[ndpath] (s0) at (0,0)      {$s^0$};
        \node[ndpath] (s1) at (2.5,0)    {$s^1$};
        \node[ndopen] (nw) at (5.0,0)    {};
        \node[nd]     (A)  at (2.5,-1.5) {};
        \node[nd]     (B)  at (5.0,-1.5) {};

        \draw[edsel] (s0) -- node[prsel] {$p{=}0.75$} (s1);
        \draw[edsel] (s1) -- node[prsel] {$p{=}0.6$}  (nw);
        \draw[ed]    (s0) -- node[pr,pos=0.55] {$p{=}0.25$} (A);
        \draw[ed]    (s1) -- node[pr,pos=0.55] {$p{=}0.4$}  (B);

        \node[st,above=1.5mm of s0] {$N{=}6$};
        \node[st,above=1.5mm of s1] {$\bar{Q}{=}0.5$,\ $N{=}2$\\pUCT${}=\mathbf{1.27}$};
        \node[st,above=1.5mm of nw] {pUCT${}=\mathbf{1.06}$};
        \node[st,below=1.5mm of A]  {$\bar{Q}{=}1.0$,\ $N{=}4$\\pUCT${}=1.15$};
        \node[st,below=1.5mm of B]  {$\bar{Q}{=}0.0$,\ $N{=}2$\\pUCT${}=0.24$};
      \end{tikzpicture}
    \end{minipage}\hfill
    \begin{minipage}[c]{0.36\textwidth}
      \caption{Nodes are selected by pUCT (\cref{eq:ez-puct}).
        Despite action~1 having twice the action value of action~0,
        action~0 is chosen because of the higher prior and lower visit count.}
      \label{fig:ez-mcts-selection}
    \end{minipage}
  \end{subfigure}

  \par\vspace{3.5ex}
  % ===== (b) Expansion and evaluation ======================================
  \begin{subfigure}[b]{\textwidth}
    \begin{minipage}[c]{0.60\textwidth}
      \centering
      \begin{tikzpicture}
        \path (-1.1,0) (7.5,0);
        \node[ndpath] (s0) at (0,0)      {$s^0$};
        \node[ndpath] (s1) at (2.5,0)    {$s^1$};
        \node[ndnew]  (s2) at (5.0,0)    {$s^2$};
        \node[nd]     (A)  at (2.5,-1.5) {};
        \node[nd]     (B)  at (5.0,-1.5) {};
        \node[ndopen] (c0) at (7.0, 0.9) {};
        \node[ndopen] (c1) at (7.0,-0.9) {};

        \draw[edsel] (s0) -- (s1);
        \draw[ed]    (s0) -- (A);
        \draw[ed]    (s1) -- (B);
        \draw[edsel] (s1) -- node[net] {$g$} (s2);
        \draw[ed]    (s2) -- node[pr,pos=0.6] {$0.35$} (c0);
        \draw[ed]    (s2) -- node[pr,pos=0.6] {$0.65$} (c1);

        \node[net] (V) at (5.0,1.35) {$v$};
        \draw[ed] (s2) -- (V);
        \node[st,right=1.5mm of V] {$v^2{=}0.6$};
      \end{tikzpicture}
    \end{minipage}\hfill
    \begin{minipage}[c]{0.36\textwidth}
      % The \todo that used to sit here moved next to %!TEX root = ../thesis.tex
% One MCTS simulation for \cref{sec:ez-mcts}, traced through the three phases.
% Worked example: |A| = 2 and gamma = 1 for legibility. The tree's Q values span
% exactly [0,1], so the tree-wide normalization of \cref{eq:ez-qnorm} is the
% identity here and the displayed values are both Q and \bar{Q}.
%
% Layout: the three phases are stacked as full-width rows rather than set side by
% side, and every tree grows LEFT TO RIGHT. Three 0.32\textwidth columns left only
% ~4.6cm per panel, which forced \tiny labels; one full-width row per phase gives
% ~8.7cm of drawing width and lets the statistics be set in \scriptsize. The
% subcaption sits beside its tree, not underneath it, so the rows stay short.
% Layout rule, to keep the panels collision-free: the simulation path runs
% straight right along y=0 with siblings on the row below; statistics of path
% nodes are therefore always set ABOVE, statistics of siblings always BELOW, and
% edge labels sit on the edge itself with a white background.
% Every panel opens with the same invisible \path at y=0, which pads all three
% tikzpictures to the same horizontal extent so s^0, s^1, s^2 line up in a column
% across the rows. It deliberately does not fix the vertical extent: (b) and (c)
% carry no statistics under their sibling nodes and would otherwise reserve the
% empty space that only (a) needs.
\begin{figure}[htbp]
  \centering
  % The subcaptions sit in a 0.36\textwidth column, too narrow to justify without
  % gaping word spaces -- set them ragged right for this figure only.
  \captionsetup[sub]{justification=raggedright, singlelinecheck=false}
  \tikzset{
    >={Stealth[round]},
    nd/.style      = {draw, circle, fill=black!8, minimum size=8mm,
                      inner sep=0pt, font=\small},
    ndpath/.style  = {nd, fill=blue!15, very thick},
    ndnew/.style   = {nd, fill=blue!15, very thick, draw=blue!60!black},
    ndopen/.style  = {draw, circle, densely dashed, black!50, minimum size=8mm,
                      inner sep=0pt, font=\small},
    net/.style     = {draw, rounded corners=2pt, fill=white, font=\scriptsize,
                      inner sep=2pt},
    ed/.style      = {->, black!55},
    edsel/.style   = {->, very thick, blue!60!black},
    edup/.style    = {->, very thick, densely dashed, black!75},
    pr/.style      = {font=\scriptsize, inner sep=1pt, black!55, fill=white},
    prsel/.style   = {font=\scriptsize, inner sep=1pt, fill=white},
    st/.style      = {font=\scriptsize, inner sep=1pt, align=center},
  }

  % ===== (a) Selection =====================================================
  \begin{subfigure}[b]{\textwidth}
    \begin{minipage}[c]{0.60\textwidth}
      \centering
      \begin{tikzpicture}
        \path (-1.1,0) (7.5,0);
        \node[ndpath] (s0) at (0,0)      {$s^0$};
        \node[ndpath] (s1) at (2.5,0)    {$s^1$};
        \node[ndopen] (nw) at (5.0,0)    {};
        \node[nd]     (A)  at (2.5,-1.5) {};
        \node[nd]     (B)  at (5.0,-1.5) {};

        \draw[edsel] (s0) -- node[prsel] {$p{=}0.75$} (s1);
        \draw[edsel] (s1) -- node[prsel] {$p{=}0.6$}  (nw);
        \draw[ed]    (s0) -- node[pr,pos=0.55] {$p{=}0.25$} (A);
        \draw[ed]    (s1) -- node[pr,pos=0.55] {$p{=}0.4$}  (B);

        \node[st,above=1.5mm of s0] {$N{=}6$};
        \node[st,above=1.5mm of s1] {$\bar{Q}{=}0.5$,\ $N{=}2$\\pUCT${}=\mathbf{1.27}$};
        \node[st,above=1.5mm of nw] {pUCT${}=\mathbf{1.06}$};
        \node[st,below=1.5mm of A]  {$\bar{Q}{=}1.0$,\ $N{=}4$\\pUCT${}=1.15$};
        \node[st,below=1.5mm of B]  {$\bar{Q}{=}0.0$,\ $N{=}2$\\pUCT${}=0.24$};
      \end{tikzpicture}
    \end{minipage}\hfill
    \begin{minipage}[c]{0.36\textwidth}
      \caption{Nodes are selected by pUCT (\cref{eq:ez-puct}).
        Despite action~1 having twice the action value of action~0,
        action~0 is chosen because of the higher prior and lower visit count.}
      \label{fig:ez-mcts-selection}
    \end{minipage}
  \end{subfigure}

  \par\vspace{3.5ex}
  % ===== (b) Expansion and evaluation ======================================
  \begin{subfigure}[b]{\textwidth}
    \begin{minipage}[c]{0.60\textwidth}
      \centering
      \begin{tikzpicture}
        \path (-1.1,0) (7.5,0);
        \node[ndpath] (s0) at (0,0)      {$s^0$};
        \node[ndpath] (s1) at (2.5,0)    {$s^1$};
        \node[ndnew]  (s2) at (5.0,0)    {$s^2$};
        \node[nd]     (A)  at (2.5,-1.5) {};
        \node[nd]     (B)  at (5.0,-1.5) {};
        \node[ndopen] (c0) at (7.0, 0.9) {};
        \node[ndopen] (c1) at (7.0,-0.9) {};

        \draw[edsel] (s0) -- (s1);
        \draw[ed]    (s0) -- (A);
        \draw[ed]    (s1) -- (B);
        \draw[edsel] (s1) -- node[net] {$g$} (s2);
        \draw[ed]    (s2) -- node[pr,pos=0.6] {$0.35$} (c0);
        \draw[ed]    (s2) -- node[pr,pos=0.6] {$0.65$} (c1);

        \node[net] (V) at (5.0,1.35) {$v$};
        \draw[ed] (s2) -- (V);
        \node[st,right=1.5mm of V] {$v^2{=}0.6$};
      \end{tikzpicture}
    \end{minipage}\hfill
    \begin{minipage}[c]{0.36\textwidth}
      % The \todo that used to sit here moved next to %!TEX root = ../thesis.tex
% One MCTS simulation for \cref{sec:ez-mcts}, traced through the three phases.
% Worked example: |A| = 2 and gamma = 1 for legibility. The tree's Q values span
% exactly [0,1], so the tree-wide normalization of \cref{eq:ez-qnorm} is the
% identity here and the displayed values are both Q and \bar{Q}.
%
% Layout: the three phases are stacked as full-width rows rather than set side by
% side, and every tree grows LEFT TO RIGHT. Three 0.32\textwidth columns left only
% ~4.6cm per panel, which forced \tiny labels; one full-width row per phase gives
% ~8.7cm of drawing width and lets the statistics be set in \scriptsize. The
% subcaption sits beside its tree, not underneath it, so the rows stay short.
% Layout rule, to keep the panels collision-free: the simulation path runs
% straight right along y=0 with siblings on the row below; statistics of path
% nodes are therefore always set ABOVE, statistics of siblings always BELOW, and
% edge labels sit on the edge itself with a white background.
% Every panel opens with the same invisible \path at y=0, which pads all three
% tikzpictures to the same horizontal extent so s^0, s^1, s^2 line up in a column
% across the rows. It deliberately does not fix the vertical extent: (b) and (c)
% carry no statistics under their sibling nodes and would otherwise reserve the
% empty space that only (a) needs.
\begin{figure}[htbp]
  \centering
  % The subcaptions sit in a 0.36\textwidth column, too narrow to justify without
  % gaping word spaces -- set them ragged right for this figure only.
  \captionsetup[sub]{justification=raggedright, singlelinecheck=false}
  \tikzset{
    >={Stealth[round]},
    nd/.style      = {draw, circle, fill=black!8, minimum size=8mm,
                      inner sep=0pt, font=\small},
    ndpath/.style  = {nd, fill=blue!15, very thick},
    ndnew/.style   = {nd, fill=blue!15, very thick, draw=blue!60!black},
    ndopen/.style  = {draw, circle, densely dashed, black!50, minimum size=8mm,
                      inner sep=0pt, font=\small},
    net/.style     = {draw, rounded corners=2pt, fill=white, font=\scriptsize,
                      inner sep=2pt},
    ed/.style      = {->, black!55},
    edsel/.style   = {->, very thick, blue!60!black},
    edup/.style    = {->, very thick, densely dashed, black!75},
    pr/.style      = {font=\scriptsize, inner sep=1pt, black!55, fill=white},
    prsel/.style   = {font=\scriptsize, inner sep=1pt, fill=white},
    st/.style      = {font=\scriptsize, inner sep=1pt, align=center},
  }

  % ===== (a) Selection =====================================================
  \begin{subfigure}[b]{\textwidth}
    \begin{minipage}[c]{0.60\textwidth}
      \centering
      \begin{tikzpicture}
        \path (-1.1,0) (7.5,0);
        \node[ndpath] (s0) at (0,0)      {$s^0$};
        \node[ndpath] (s1) at (2.5,0)    {$s^1$};
        \node[ndopen] (nw) at (5.0,0)    {};
        \node[nd]     (A)  at (2.5,-1.5) {};
        \node[nd]     (B)  at (5.0,-1.5) {};

        \draw[edsel] (s0) -- node[prsel] {$p{=}0.75$} (s1);
        \draw[edsel] (s1) -- node[prsel] {$p{=}0.6$}  (nw);
        \draw[ed]    (s0) -- node[pr,pos=0.55] {$p{=}0.25$} (A);
        \draw[ed]    (s1) -- node[pr,pos=0.55] {$p{=}0.4$}  (B);

        \node[st,above=1.5mm of s0] {$N{=}6$};
        \node[st,above=1.5mm of s1] {$\bar{Q}{=}0.5$,\ $N{=}2$\\pUCT${}=\mathbf{1.27}$};
        \node[st,above=1.5mm of nw] {pUCT${}=\mathbf{1.06}$};
        \node[st,below=1.5mm of A]  {$\bar{Q}{=}1.0$,\ $N{=}4$\\pUCT${}=1.15$};
        \node[st,below=1.5mm of B]  {$\bar{Q}{=}0.0$,\ $N{=}2$\\pUCT${}=0.24$};
      \end{tikzpicture}
    \end{minipage}\hfill
    \begin{minipage}[c]{0.36\textwidth}
      \caption{Nodes are selected by pUCT (\cref{eq:ez-puct}).
        Despite action~1 having twice the action value of action~0,
        action~0 is chosen because of the higher prior and lower visit count.}
      \label{fig:ez-mcts-selection}
    \end{minipage}
  \end{subfigure}

  \par\vspace{3.5ex}
  % ===== (b) Expansion and evaluation ======================================
  \begin{subfigure}[b]{\textwidth}
    \begin{minipage}[c]{0.60\textwidth}
      \centering
      \begin{tikzpicture}
        \path (-1.1,0) (7.5,0);
        \node[ndpath] (s0) at (0,0)      {$s^0$};
        \node[ndpath] (s1) at (2.5,0)    {$s^1$};
        \node[ndnew]  (s2) at (5.0,0)    {$s^2$};
        \node[nd]     (A)  at (2.5,-1.5) {};
        \node[nd]     (B)  at (5.0,-1.5) {};
        \node[ndopen] (c0) at (7.0, 0.9) {};
        \node[ndopen] (c1) at (7.0,-0.9) {};

        \draw[edsel] (s0) -- (s1);
        \draw[ed]    (s0) -- (A);
        \draw[ed]    (s1) -- (B);
        \draw[edsel] (s1) -- node[net] {$g$} (s2);
        \draw[ed]    (s2) -- node[pr,pos=0.6] {$0.35$} (c0);
        \draw[ed]    (s2) -- node[pr,pos=0.6] {$0.65$} (c1);

        \node[net] (V) at (5.0,1.35) {$v$};
        \draw[ed] (s2) -- (V);
        \node[st,right=1.5mm of V] {$v^2{=}0.6$};
      \end{tikzpicture}
    \end{minipage}\hfill
    \begin{minipage}[c]{0.36\textwidth}
      % The \todo that used to sit here moved next to \input{figures/ez-mcts-simulation}
      % in sections/efficientzero.tex: \todo uses \marginpar, and a \marginpar inside
      % a subfigure (a minipage) makes LaTeX abort with ``Float(s) lost''.
      \caption{The selected node is expanded. The latent state is rolled out once
        via $g$ and value statistics are initialized.}
      \label{fig:ez-mcts-expansion}
    \end{minipage}
  \end{subfigure}

  \par\vspace{3.5ex}
  % ===== (c) Backup ========================================================
  \begin{subfigure}[b]{\textwidth}
    \begin{minipage}[c]{0.60\textwidth}
      \centering
      \begin{tikzpicture}
        \path (-1.1,0) (7.5,0);
        \node[ndpath] (s0) at (0,0)      {$s^0$};
        \node[ndpath] (s1) at (2.5,0)    {$s^1$};
        \node[ndnew]  (s2) at (5.0,0)    {$s^2$};
        \node[nd]     (A)  at (2.5,-1.5) {};
        \node[nd]     (B)  at (5.0,-1.5) {};

        \draw[ed] (s0) -- (A);
        \draw[ed] (s1) -- (B);
        \draw[edup] (s1) -- node[prsel] {$r{=}0.2$} (s0);
        \draw[edup] (s2) -- node[prsel] {$r{=}0.4$} (s1);

        \node[st,above=1.5mm of s0] {$G^0{=}1.2$\\$N{:}\,6{\to}7$\\$Q_\Sigma{:}\,3.6{\to}4.8$};
        \node[st,above=1.5mm of s1] {$G^1{=}1.0$\\$N{:}\,2{\to}3$\\$Q_\Sigma{:}\,0.6{\to}1.6$};
        \node[st,above=1.5mm of s2] {$G^2{=}v^2{=}0.6$\\$N{:}\,0{\to}1$\\$Q_\Sigma{:}\,0{\to}0.6$};
      \end{tikzpicture}
    \end{minipage}\hfill
    \begin{minipage}[c]{0.36\textwidth}
      \caption{The leaf estimate is backed up through the path. Each node accumulates
        its return and the visit count is incremented.}
      \label{fig:ez-mcts-backup}
    \end{minipage}
  \end{subfigure}

  \caption[One MCTS simulation]{A single iteration of Monte-Carlo
    tree search, where $\lvert\mathcal{A}\rvert=2$ and $\gamma=1$.
    Search starts from a true root latent $s^0$ and is then continued by imagined rollouts
    to imagined latents $s^1$ and $s^2$.
    $Q_\Sigma$ represents the sum of all backed-up Q values, $N$ represents the visit count,
    $\bar{Q}=\frac{Q_\Sigma}{N}$ the average Q value of the node, and $\text{pUCT}$ the actual score
    according to the pUCT rule~(\cref{eq:ez-puct}).
    Colored nodes and edges indicate the search path while the edges are labeled with the prior $p$.}
  \label{fig:ez-mcts-simulation}
\end{figure}

      % in sections/efficientzero.tex: \todo uses \marginpar, and a \marginpar inside
      % a subfigure (a minipage) makes LaTeX abort with ``Float(s) lost''.
      \caption{The selected node is expanded. The latent state is rolled out once
        via $g$ and value statistics are initialized.}
      \label{fig:ez-mcts-expansion}
    \end{minipage}
  \end{subfigure}

  \par\vspace{3.5ex}
  % ===== (c) Backup ========================================================
  \begin{subfigure}[b]{\textwidth}
    \begin{minipage}[c]{0.60\textwidth}
      \centering
      \begin{tikzpicture}
        \path (-1.1,0) (7.5,0);
        \node[ndpath] (s0) at (0,0)      {$s^0$};
        \node[ndpath] (s1) at (2.5,0)    {$s^1$};
        \node[ndnew]  (s2) at (5.0,0)    {$s^2$};
        \node[nd]     (A)  at (2.5,-1.5) {};
        \node[nd]     (B)  at (5.0,-1.5) {};

        \draw[ed] (s0) -- (A);
        \draw[ed] (s1) -- (B);
        \draw[edup] (s1) -- node[prsel] {$r{=}0.2$} (s0);
        \draw[edup] (s2) -- node[prsel] {$r{=}0.4$} (s1);

        \node[st,above=1.5mm of s0] {$G^0{=}1.2$\\$N{:}\,6{\to}7$\\$Q_\Sigma{:}\,3.6{\to}4.8$};
        \node[st,above=1.5mm of s1] {$G^1{=}1.0$\\$N{:}\,2{\to}3$\\$Q_\Sigma{:}\,0.6{\to}1.6$};
        \node[st,above=1.5mm of s2] {$G^2{=}v^2{=}0.6$\\$N{:}\,0{\to}1$\\$Q_\Sigma{:}\,0{\to}0.6$};
      \end{tikzpicture}
    \end{minipage}\hfill
    \begin{minipage}[c]{0.36\textwidth}
      \caption{The leaf estimate is backed up through the path. Each node accumulates
        its return and the visit count is incremented.}
      \label{fig:ez-mcts-backup}
    \end{minipage}
  \end{subfigure}

  \caption[One MCTS simulation]{A single iteration of Monte-Carlo
    tree search, where $\lvert\mathcal{A}\rvert=2$ and $\gamma=1$.
    Search starts from a true root latent $s^0$ and is then continued by imagined rollouts
    to imagined latents $s^1$ and $s^2$.
    $Q_\Sigma$ represents the sum of all backed-up Q values, $N$ represents the visit count,
    $\bar{Q}=\frac{Q_\Sigma}{N}$ the average Q value of the node, and $\text{pUCT}$ the actual score
    according to the pUCT rule~(\cref{eq:ez-puct}).
    Colored nodes and edges indicate the search path while the edges are labeled with the prior $p$.}
  \label{fig:ez-mcts-simulation}
\end{figure}

      % in sections/efficientzero.tex: \todo uses \marginpar, and a \marginpar inside
      % a subfigure (a minipage) makes LaTeX abort with ``Float(s) lost''.
      \caption{The selected node is expanded. The latent state is rolled out once
        via $g$ and value statistics are initialized.}
      \label{fig:ez-mcts-expansion}
    \end{minipage}
  \end{subfigure}

  \par\vspace{3.5ex}
  % ===== (c) Backup ========================================================
  \begin{subfigure}[b]{\textwidth}
    \begin{minipage}[c]{0.60\textwidth}
      \centering
      \begin{tikzpicture}
        \path (-1.1,0) (7.5,0);
        \node[ndpath] (s0) at (0,0)      {$s^0$};
        \node[ndpath] (s1) at (2.5,0)    {$s^1$};
        \node[ndnew]  (s2) at (5.0,0)    {$s^2$};
        \node[nd]     (A)  at (2.5,-1.5) {};
        \node[nd]     (B)  at (5.0,-1.5) {};

        \draw[ed] (s0) -- (A);
        \draw[ed] (s1) -- (B);
        \draw[edup] (s1) -- node[prsel] {$r{=}0.2$} (s0);
        \draw[edup] (s2) -- node[prsel] {$r{=}0.4$} (s1);

        \node[st,above=1.5mm of s0] {$G^0{=}1.2$\\$N{:}\,6{\to}7$\\$Q_\Sigma{:}\,3.6{\to}4.8$};
        \node[st,above=1.5mm of s1] {$G^1{=}1.0$\\$N{:}\,2{\to}3$\\$Q_\Sigma{:}\,0.6{\to}1.6$};
        \node[st,above=1.5mm of s2] {$G^2{=}v^2{=}0.6$\\$N{:}\,0{\to}1$\\$Q_\Sigma{:}\,0{\to}0.6$};
      \end{tikzpicture}
    \end{minipage}\hfill
    \begin{minipage}[c]{0.36\textwidth}
      \caption{The leaf estimate is backed up through the path. Each node accumulates
        its return and the visit count is incremented.}
      \label{fig:ez-mcts-backup}
    \end{minipage}
  \end{subfigure}

  \caption[One MCTS simulation]{A single iteration of Monte-Carlo
    tree search, where $\lvert\mathcal{A}\rvert=2$ and $\gamma=1$.
    Search starts from a true root latent $s^0$ and is then continued by imagined rollouts
    to imagined latents $s^1$ and $s^2$.
    $Q_\Sigma$ represents the sum of all backed-up Q values, $N$ represents the visit count,
    $\bar{Q}=\frac{Q_\Sigma}{N}$ the average Q value of the node, and $\text{pUCT}$ the actual score
    according to the pUCT rule~(\cref{eq:ez-puct}).
    Colored nodes and edges indicate the search path while the edges are labeled with the prior $p$.}
  \label{fig:ez-mcts-simulation}
\end{figure}

\todo{how are value statistics initialized? (panel (b), \cref{fig:ez-mcts-expansion})}

\paragraph{Selection}\label{sec:ez-mcts-selection}
Starting at a given root state $s^0$, an action is chosen until a yet unexpanded
node is reached.
Actions are chosen using the pUCT rule~\cite{schrittwieser_mastering_2020}:

\begin{equation}\label{eq:ez-puct}
  a^k = \argmax_a\left[
    \bar{Q}(s,a)
    + p_a(s)
      \cdot\frac{\sqrt{\sum_b N(s,b)}}{1+N(s,a)}
      \cdot\left(c_1+\log\frac{\sum_b N(s,b)+c_2+1}{c_2}\right)
  \right],
\end{equation}

where constants $c_1,c_2$ control the exploration term, $N(s,a)$ is the number
of times action $a$ has already been selected in state $s$, $p_a(s)$ is the
policy prior, and $\bar{Q}(s,a)$ is a normalized, empirical action value.
Writing $s'=g(s,a;\theta)$ for the child reached by $a$ and $r(s,a)$ for the
reward on that edge:

\begin{equation}\label{eq:ez-qnorm}
  \begin{aligned}
    \bar{Q}(s,a)
      &= \frac{Q(s,a)-\min_{(s',a')\in\mathit{Tree}} Q(s',a')}
              {\max_{(s',a')\in\mathit{Tree}} Q(s',a')-\min_{(s',a')\in\mathit{Tree}} Q(s',a')}, \\
    Q(s,a) &= r(s,a)+\gamma\,\bar{v}(s'), \\
    \bar{v}(s')  &= \frac{Q_\Sigma(s')}{N(s')}.
  \end{aligned}
\end{equation}

% TODO: $r(s,a)$ is not a prediction of the model: EfficientZero has no per-step
% reward head and recovers the edge reward as a difference of consecutive value
% prefixes (\cref{sec:ez-value-prefix}). Say so here --- as written, the reader
% meets $r$ in \cref{eq:ez-qnorm} and \cref{eq:ez-backup} and only learns 150
% lines later that nothing produces it directly.

To improve exploration, the root priors are perturbed with Dirichlet noise
before the search begins.

\begin{equation}\label{eq:ez-dirichlet}
  p_a(s^0) \leftarrow (1-\rho)\,p_a(s^0)+\rho\,\eta_a,
  \qquad \eta\sim\operatorname{Dir}(\xi),
\end{equation}

where $\rho\in[0,1]$ interpolates between the prior and Dirichlet noise and $\xi$
controls its shape.

\paragraph{Expansion and Evaluation}\label{sec:ez-mcts-expansion}

When a sequence of actions has been selected and an unexpanded node is reached,
it is expanded: the previous latent state is unrolled one step to get the new
latent $s^{k+1}=g(s^k,a^k;\theta)$ and the value statistics are initialized.

\paragraph{Backup}\label{sec:ez-mcts-backup}

The leaf value is propagated back along the search path
$s^0, a^1, s^1, \dots, a^l, s^l$.
For every node on the path, the backed-up return is the discounted sum of the
rewards below it, bootstrapped with the leaf estimate:

\begin{equation}\label{eq:ez-backup}
  G^k=\sum_{\tau=0}^{l-1-k}\gamma^{\tau}\,r^{k+1+\tau}\;+\;\gamma^{l-k}\,v^l,
\end{equation}

which the backup from the leaf computes recursively as $G^l=v^l$ and
$G^k=r^{k+1}+\gamma\,G^{k+1}$.
Each node on the path then accumulates its return and one visit:

\begin{equation}\label{eq:ez-backup-update}
  Q_\Sigma(s^k) \leftarrow Q_\Sigma(s^k)+G^k,
  \qquad
  N(s^k) \leftarrow N(s^k)+1,
  \qquad k=l,\dots,0.
\end{equation}

\section{The Self-Play and Training Loop}\label{sec:ez-loop}

EfficientZero uses a parallel architecture with multiple worker types.
\Cref{fig:ez-workers} shows how they are wired together.

%!TEX root = ../thesis.tex
% Worker architecture for \cref{sec:ez-loop}: the four worker types, the replay
% buffer they share, and the parameter feedback that closes the loop.
% Layout rule: the five stages form a rectangular cycle, traversed clockwise
% --- right along the top row, down the right edge, left along the bottom row,
% and up the left edge, where the parameters close it. The environment hangs
% off the cycle above the data workers, since it is not part of the algorithm.
% Every edge is horizontal or vertical: keep it that way. An earlier draft put
% the GPU workers under the environment instead of under the data workers,
% which forced the parameter arrow into a diagonal that crossed the $a_t$ edge.
% Spacing constraint: an edge label sits in the gap between two boxes, so it
% must be NARROWER than that gap or it slides under them. Columns are 5.2cm
% apart and boxes are 3cm wide, leaving a 2.2cm (62.6pt) gap; at \scriptsize
% "prioritized positions" alone is 74pt, which is why such labels are stacked
% over two lines. Widen the columns before un-stacking them. "refreshed
% targets" (61.7pt) is single-line because it spans the whole bottom row, not a
% gap, and so has 70pt of clearance on either side.
% Solid arrows carry data, dashed arrows carry parameters.
\begin{figure}[htbp]
  \centering
  \begin{tikzpicture}[
      >={Stealth[round]}, font=\small,
      worker/.style = {draw, rounded corners=2pt, fill=blue!8,
                       minimum height=9mm, minimum width=30mm, align=center},
      store/.style  = {draw, fill=black!8,
                       minimum height=9mm, minimum width=30mm, align=center},
      env/.style    = {draw, densely dashed, fill=white,
                       minimum height=9mm, minimum width=30mm, align=center},
      flow/.style   = {->, thick, black!70},
      param/.style  = {->, thick, black!70, densely dashed},
      lbl/.style    = {font=\scriptsize, inner sep=2.5pt, align=center},
    ]
    % --- the environment, outside the cycle ----------------------------------
    \node[env]    (env)  at (0,2.3)     {Environment};

    % --- top row: acting -----------------------------------------------------
    \node[worker] (data) at (0,0)       {Data workers};
    \node[store]  (buf)  at (5.2,0)     {Replay buffer};
    \node[worker] (roll) at (10.4,0)    {Rollout workers};

    % --- bottom row: learning ------------------------------------------------
    \node[worker] (gpu)  at (0,-3.1)    {GPU workers};
    \node[worker] (rean) at (10.4,-3.1) {Reanalyze workers};

    % --- environment interaction, as a parallel pair -------------------------
    \draw[flow] ([xshift=-8mm]env.south)  --
                node[lbl,left]  {$o_t,\,u_t$} ([xshift=-8mm]data.north);
    \draw[flow] ([xshift=8mm]data.north)  --
                node[lbl,right] {$a_t$}       ([xshift=8mm]env.south);

    % --- data path: clockwise around the cycle -------------------------------
    \draw[flow] (data) -- node[lbl,above] {$(o,u,\pi_t,\nu_t)$} (buf);
    \draw[flow] (buf)  -- node[lbl,above] {prioritized\\positions} (roll);
    \draw[flow] (roll) -- node[lbl,left]  {training\\batches} (rean);
    \draw[flow] (rean) -- node[lbl,above] {refreshed targets} (gpu);

    % --- parameter feedback: closes the cycle, and reaches the searchers -----
    \draw[param] (gpu) -- node[lbl,left] {$\theta$} (data);
    \draw[param,rounded corners=3pt]
      (gpu.south) -- (0,-4.1) -- node[lbl,fill=white] {$\theta$} (10.4,-4.1)
                  -- (rean.south);
  \end{tikzpicture}
  \caption[The EfficientZero worker architecture]{%
    The parallel worker architecture of EfficientZero.
    Data workers generate fresh trajectories and store them inside the replay buffer.
    Rollout and reanalyze workers sample from the replay buffer and prepare data for training.
    GPU workers calculate losses and update the model parameters.
  }
  \label{fig:ez-workers}
\end{figure}


\paragraph{Data workers}

New data is generated by interacting with the environment.
At every time step $t$, a worker encodes the latest observation into
$s_t = h(o_t;\theta)$ and runs MCTS from this root.
The resulting root visit count distribution is then sampled using a temperature
$T$ and the action passed to the environment:

\begin{equation}\label{eq:ez-visit-policy}
  a_t\sim\pi_t(a)=\frac{N(a)^{1/T}}{\sum_b N(b)^{1/T}},
\end{equation}

where $N(a)$ is the number of times the search selected action $a$ at the
root.
The temperature $T$ interpolates between sampling proportionally to the visit
counts ($T=1$) and greedily playing the most-visited action
($T\rightarrow 0$).

The resulting observation-reward tuples $(o,u)$ are stored in a replay buffer,
together with the root visit count distribution and root state value.
As soon as an episode ends, the environment state is reset and data collection starts anew.

\paragraph{Rollout workers}

The rollout workers repeatedly sample positions from the replay buffer and
prepare training batches.
Each trajectory position inside the replay buffer has a sampling priority associated with it~\cite{schaul_prioritized_2016}.
This priority is proportional to the loss of the value network $v$ evaluated at this position.
Prioritizing such positions with a high loss accelerates learning.

\paragraph{Reanalyze workers}

Prepared batches then get processed by reanalyze workers.
\emph{MuZero Reanalyze} is an extension to the basic MuZero algorithm that improves the freshness of
data by re-running MCTS on old data with updated model weights, contributing to data reuse and sample efficiency.
Reanalyzed batches are finally passed to the GPU workers where the model is trained.

This architecture is an instance of generalized policy iteration:
MCTS serves as the policy improvement operator~\cite{grill_monte-carlo_2020}, while learning from stored experience
improves the value estimates.

\section{Training Objective}\label{sec:ez-loss}

The reanalyzed batches are then used to jointly train all networks.
A batch entry starting at position $t$ is encoded once as $s^0=h(o_t;\theta)$ and then
unrolled for $K$ steps with the actions that were actually taken,
$s^k=g(s^{k-1},a_{t+k};\theta)$.
Each unrolled latent $s^k$ is supervised with the targets located at the real
position $t+k$. These include the actual reward and actual observation as well
as MCTS root visit distribution and root value target.

\paragraph{Policy target}
The policy head is trained toward the MCTS root visit distribution $\pi_{t+k}$
to align the current policy with the improved MCTS policy.

\paragraph{Value target}
The value head is trained toward an $n$-step return of actual rewards,
bootstrapped with the root value that MCTS produced:

\begin{equation}\label{eq:ez-value-target}
  z_t=\sum_{i=0}^{n-1}\gamma^{i}\,u_{t+i+1}+\gamma^{n}\,v^{\text{MCTS}}_{t+n},
\end{equation}

where $u$ are the actual rewards observed by the data workers and
$v^{\text{MCTS}}_{t+n}$ is the MCTS root value.

\paragraph{Value-prefix target}
The value-prefix head is trained toward $x_{t,k}$ of \cref{eq:ez-value-prefix},
the discounted reward accumulated between the root and unroll step $k$.
It replaces MuZero's per-step reward target.

\paragraph{Consistency target}
The dynamics network is additionally supervised by a consistency loss
calculated from the predicted successor latent $s^k=g(s^{k-1},a_{t+k};\theta)$
and true successor latent $h(o_{t+k};\theta)$.

\paragraph{Total loss}
The four terms are weighted and summed over the unroll horizon $0\dots K$:

\begin{equation}\label{eq:ez-loss}
  \begin{split}
    \mathcal{L}(\theta)=\sum_{k=0}^{K}\Big[
      &\lambda_p\,\ell\!\left(\pi_{t+k},p(s^k;\theta)\right)
       + \lambda_v\,\ell\!\left(z_{t+k},v(s^k;\theta)\right)\\
      &+ \lambda_x\,\ell\!\left(x_{t,k},\vpnet(s^0,a_{t+1:t+k};\theta)\right)
       + \lambda_c\,\mathcal{L}^k_{\text{consistency}}\Big],
  \end{split}
\end{equation}

where $\ell$ is the cross entropy between a target and the categorical
prediction of the corresponding head.
The value loss $\lambda_v\,\ell\!\left(z_{t+k},v(s^k;\theta)\right)$ is used to update
the replay buffer priorities.

\section{EfficientZero's Improvements}\label{sec:ez-improvements}

The EfficientZero paper~\cite{ye_mastering_2021} identifies three weaknesses of
the vanilla architecture. To improve the learner performance, they introduce
three changes:

\subsection{Self-Supervised Consistency Loss}\label{sec:ez-consistency}

In MuZero, training the dynamics function happens implicitly through the prediction
of rewards, policies and values.
In a reward-sparse environment, this signal is often not strong enough and causes
the dynamics function $g$ to lose its predictive power.

EfficientZero adds an additional, SimSiam-inspired~\cite{chen_exploring_2020}
training signal.
The idea is that the imagined next state $s^{k}=g(s^{k-1},a^{k};\theta)$
and the encoded real next state $h(o_{t+k};\theta)$ should be similar.
To this end, both states are passed through a projection network $\projnet$, the
imagined branch additionally through a prediction head $\prednet$, and the two
outputs are pulled together with a negative cosine similarity loss:

\begin{equation}\label{eq:ez-consistency}
  \mathcal{L}_{\text{consistency}}
  = -\operatorname{cossim}\!\Big(
      \operatorname{sg}\big(\projnet(h(o_{t+k};\theta);\theta)\big),\;
      \prednet\big(\projnet(s^{k};\theta);\theta\big)
    \Big),
\end{equation}

where the stop-gradient operator $\operatorname{sg}$ prevents the trivial
solution of mapping every observation to a constant.
The loss is applied at every unroll step, giving the dynamics function a dense
training signal.
The idea is closely related to SPR~\cite{schwarzer_data-efficient_2021}, which
uses the momentum-encoder variant BYOL~\cite{grill_bootstrap_2020} for the
same purpose in a model-free agent.

\subsection{End-to-End Value-Prefix Prediction}\label{sec:ez-value-prefix}

MuZero's dynamics function must predict the reward of every single imagined
step.
Predicting exactly when a reward is given is a hard problem.
For each additional imagined unroll step, error terms compound.
EfficientZero sidesteps this issue by directly predicting the \emph{value prefix} $x_{t,k}$,
a sum of discounted rewards over a given sequence of actions:

\begin{equation}\label{eq:ez-value-prefix}
  x_{t,k}=\sum_{i=0}^{k-1}\gamma^{i}\,u_{t+i+1}.
\end{equation}

When the search requires an individual edge reward, it is recovered as the
difference of consecutive prefixes.

\subsection{Model-Based Off-Policy Correction}\label{sec:ez-off-policy}

The value target of \cref{eq:ez-value-target} sums $n$ rewards that were
collected by a potentially older policy and might therefore be biased.
Off-policy bias of multi-step targets is a classical problem.
EfficientZero corrects the target by varying the bootstrap horizon:
a trajectory that is $\tau$ updates old uses a shortened horizon $l\leq n$,
causing less stale data to enter the bootstrapped value target.

\begin{equation}\label{eq:ez-corrected-target}
  z_t=\sum_{i=0}^{l-1}\gamma^{i}\,u_{t+i+1}+\gamma^{l}\,v^{\text{MCTS}}_{t+l}.
\end{equation}
